% appendix_data_issues.tex: the credit-bureau data: files and frequencies, sample restrictions, self-calculated measures, and the data issues, stated against Gibbs, Guttman-Kenney, Lee, Nelson, van der Klaauw and Wang (2025). Canonical list: memos/data_issues_ledger.md.
\section{Credit-Bureau Data: Construction, Measures, and Their Limits}\label{appendix:data}

The panel is the Georgia Institute of Technology Scheller College of Business consumer credit panel drawn from a national credit bureau's files, which \citet{gibbs_consumer_2025} list in their Table 4 and describe in their Online Appendix G.1. Their Section 3.1 asks every user of credit reporting data to state four things: which files the measures are calculated from, at what frequency, under which sample restrictions, and which measures the researcher calculates and by what formula. Appendix \ref{appendix:data:four} answers the four for every exhibit in the paper. Section \ref{appendix:data:issues} records each feature of the raw files that bears on a measure, what the paper does about it, and where it matters, with the section of the survey that treats it. Section \ref{appendix:data:adds} lists what the panel adds to the survey's account.

\subsection{The loan types as parallel laboratories}\label{appendix:data:loantypes}

Table \ref{tab:loan_types} sets the loan types side by side: the contract, the population and horizon, the identification each offers, and what each delivers in the paper (Section \ref{sec:autodata}).

\begin{landscape}
\begin{table}[H]
\centering
\caption{The loan types as parallel laboratories}
\label{tab:loan_types}
% replication: hand-written (Draft/tables/tab_loan_types.tex); numbers checked against results/summary_stats_families.csv
\scriptsize
% tables/tab_loan_types.tex — the loan types as parallel laboratories (hand-written; facts from the sections cited in the notes).
\begin{tabularx}{0.96\textheight}{@{}>{\raggedright\arraybackslash}p{1.7cm}>{\raggedright\arraybackslash}X>{\raggedright\arraybackslash}X>{\raggedright\arraybackslash}X>{\raggedright\arraybackslash}X@{}}
\toprule
Loan type & Contract & Population and horizon & Identification available & Delivered in this paper \\
\midrule
Auto loans & Fixed payment, 12 to 96 months, amortizing; collateral repossessed; deficiency pursuable in most states & Wide participation, including the low end of the income and score distributions that asset markets and laboratory samples miss; one to six years & Payment and term margins within lender, place and score-by-income cells; within-borrower designs; the deficiency-judgment contrast; the model's own run-off & The main panel (7.6 million loans, 15,532 lenders): the payment margin at population scale and its heterogeneity (Sections \ref{sec:mainresults}, \ref{sec:heterogeneity}); the bureau estimate of the rate (Appendix \ref{appendix:ladder}) \\
Credit cards & Revolving; the payment is an issuer-set minimum formula on the balance plus interest; no collateral & The same low end of the income distribution; months to about four years, the duration of the minimum-payment path (Appendix \ref{appendix:cards}) & Issuer-assigned changes to the minimum-payment formula and promotional-rate expirations (Appendix \ref{appendix:strategies}); the balance and formula margins & The revolving-credit extension and its feasibility in the archives (Appendix \ref{appendix:cards}) \\
First mortgages & Fixed or adjustable rate, 30 years, amortizing; the house as collateral; recourse varies by state & Homeowners; the longest horizons, up to 30 years & The equity margin from house prices by origination cohort within zip-by-quarter cells; adjustable-rate resets at contract-fixed dates; the recourse contrast; the deficiency contrast on the term margin & The equity margin and the walk-away signatures (Section \ref{sec:equity}, Appendix \ref{appendix:equity}); the reset facts (Appendix \ref{sec:armreset}); the recourse measurement of the penalty channel \\
Student loans & Statutory terms; income-driven plans with payments set by income; no collateral & Borrowers with post-secondary education; ten to 25 years & The October 2023 payment resumption, a bill on another account at a date fixed by statute (Appendix \ref{sec:studentloan}) & The response of auto default to an unenforced bill: zero at every horizon (Section \ref{sec:equity}, Appendix \ref{sec:studentloan}) \\
Personal installment; home-equity loans & Fixed payment, amortizing; unsecured (median \$2,791 over 24 months), or secured by the house & Overlapping the auto population; months to a few years, and up to 30 & The same margins as auto loans; within-borrower term margin on multi-loan borrowers & Panels constructed; within borrower the personal-loan term margin at 24 months is 0.3 in every score tercile, as for auto loans; the payment margin needs family-specific windows (Appendix \ref{appendix:strategies}) \\
\bottomrule
\end{tabularx}

\vspace{0.5em}
\begin{minipage}{0.92\textwidth}
\noindent \small \textit{Notes:} Contract features and identification strategies as used in the sections cited; the population columns describe the part of the income and credit-score distribution each loan type covers in the one-percent panel. Facts on each row are those reported in the cited sections.
\end{minipage}
\end{table}
\end{landscape}

\subsection{Files, frequencies, restrictions and formulas by exhibit}\label{appendix:data:four}

The institution holds two nested consumer panels drawn from the same bureau. The one-percent panel contains trade files for 202 archive months, semiannual from 2005 to 2009 and monthly from January 2010 to December 2025, a monthly score file for the same months, a quarterly attribute file with age and modeled income, and performance attributes from 2010. The ten-percent panel contains trade files for 57 archive months, semiannual from 2005 to 2017, quarterly from 2018 to 2022, monthly for part of 2023 and quarterly to April 2024, and a monthly score file from 2010, 195 files of about 31 million consumers each with two score columns, whose consumer keys match the trade files in full. Age and income exist only for the one-percent consumers, who are nested in the ten-percent panel. Table \ref{tab:data_four} gives, for each exhibit type, the files and frequency, the restrictions, and the self-calculated measures.

\begingroup
\footnotesize
\setstretch{1.0}
\begin{longtable}{@{}>{\raggedright\arraybackslash}p{2.5cm}>{\raggedright\arraybackslash}p{3.5cm}>{\raggedright\arraybackslash}p{3.6cm}>{\raggedright\arraybackslash}p{5.2cm}@{}}
\caption{The four disclosures of \citet[Section 3.1]{gibbs_consumer_2025} for every exhibit type}\label{tab:data_four}\\
% replication: hand-written from the cited survey; no data
\toprule
Exhibit & Files and frequency & Restrictions & Self-calculated measures \\
\midrule
\endfirsthead
\toprule
Exhibit & Files and frequency & Restrictions & Self-calculated measures \\
\midrule
\endhead
\bottomrule
\endfoot
Auto payment and term margins; exposure by group; run-off; ladder (Tables \ref{tab:main_margins} upper panel, \ref{tab:exposure}, \ref{tab:runoff_shape}, \ref{tab:bureau_ladder}, \ref{tab:bureau_rho_groups}; Figures \ref{fig:exposure_groups}, \ref{fig:exposure_geo}, \ref{fig:runoff_profiles}, \ref{fig:runoff_shape}) & One-percent trade files, semiannual 2005--09 and monthly 2010-01 to 2025-12; one-percent score file at the month of first sight; one-percent attribute file at the nearest earlier quarter within a year & Auto flag; contract length 12 to 96 months; positive scheduled payment; score in $(0, 851)$; positive income; lender key present; origination date parsed; for a window of $K$ months, contract length at least $K$ and opening at least $K$ months before the last archive & Loan record: one per consumer key, lender key and origination date, origination fields at first sight, terminal fields at last sight. Charge-off within $K$ months: a positive charge-off amount at an archive within $K$ months of the origination date. Ninety-day delinquency within $K$ months: a manner-of-payment rating of 4 or worse at such an archive. Elasticity: linear-probability coefficient on $\log m$ or $\log T$ over the window's default rate. Run-off: the same by loan-age bin among loans surviving to the bin. Early payoff: no charge-off and a closure, a zero balance or a cessation of reporting more than six months before both maturity and the last archive. Exit hazard: payoffs over years of exposure, $0.22$ per year on autos (\texttt{results/panel\_exit\_hazards\_auto.csv}) \\
Within-borrower margins; term margin by score; deficiency contrast (Tables \ref{tab:main_margins} lower panel, \ref{tab:within_by_group}, \ref{tab:term_by_score}, \ref{tab:term_deficiency}) & Ten-percent trade files, 57 archives; ten-percent score file, monthly from 2010, at first sight; age and income from the one-percent attribute file for the one-percent consumers; state deficiency rules from an external file & Auto family as above; within borrower, consumers with two or more auto loans under the same consumer key; the covariate-matched subsample is the loans with a score and an income & Charge-off as above, dated by the last-payment date on this extract; within-borrower elasticity within borrower and lender-year cells; selection correction: within-borrower minus across-borrower term elasticity \\
Distress persistence (Appendix \ref{appendix:measuredrho}) & Ten-percent trade files, consecutive archives, mean gap $4.2$ months & Trades with a rating of 30 days past due or worse, no charge-off, positive balance at the earlier archive & Outcome at the next archive: cured (rating current), charge-off (positive charge-off amount or rating charge-off or repossession), still delinquent, absent. Per-gap persistence $\rho_{\text{gap}}$: still and charge-off over observed outcomes; $\rho_{\text{annual}} = \rho_{\text{gap}}^{12/\text{gap}}$; the absent share is reported, $13.7$ percent on autos and $19.1$ on mortgages (\texttt{results/cure\_rates\_v2.csv}) \\
Other trades of auto defaulters (Table \ref{tab:liquidity}, Panel A) & Ten-percent trade files & Auto defaulters holding at least one other open trade with no charge-off at the event & Event: the first archive with a positive charge-off amount on any auto trade; status of the other trades at the archives nearest one month before and six months after, each within two months of its target; paying: a last-payment date within three months of the post archive \\
Equity margin (Table \ref{tab:equity_margin}, Appendix \ref{appendix:equity}) & One-percent mortgage panel, archives 2005--18 (semiannual to 2009, monthly from 2010), collapsed to loan-quarters; Zillow home value index by five-digit zip, monthly, averaged within the quarter & Thirty-year first mortgages originated 2000--12; an index in the origination quarter (88 percent of loans meeting the contract filters) & At risk from the first quarter in the panel to the quarter of the first 90-day delinquency, payoff or last sight; log house-price change since origination; origination rate by the annuity identity; hazard in percent per quarter; elasticity: coefficient over the base hazard \\
Adjustable-rate resets (Tables \ref{tab:arm_classes}, \ref{tab:arm_samples}; Figures \ref{fig:arm_event}, \ref{fig:arm_anticipation}) & One-percent trade files, monthly from 2010, first mortgages originated 2004--08 observed 2005--18; one-year Treasury rate from an external file & Payment changes of more than five percent at an amortizing balance; hybrid reset ages within three months of 36, 60 or 84 & Note rate from the balance path: with remaining term $n$, balance $B$ and monthly principal reduction $\pi$, $\pi/B = (r/12)/((1+r/12)^n - 1)$; principal-and-interest payment from $r$; escrow as the remainder of the reported payment; classes by the implied-rate and escrow changes over eight months on either side (Table \ref{tab:arm_classes} notes); first 90-day delinquency as the outcome \\
Student-loan resumption (Tables \ref{tab:student_takeup}, \ref{tab:student_groups}; Figure \ref{fig:student_profile}) & One-percent trade files: student trades at 2020-02, 2023-06, 2023-09 to 2024-01, 2024-10 and 2025-01; auto loan-months 2021-07 to 2025-12 & Consumers with a positive student-loan balance in 2023-09; auto loans open in 2023-01 & Exposure: the resumed scheduled payment in 2023-11 to 2024-01 in hundreds of dollars; take-up: a positive actual-payment amount on any student trade; past due: a positive past-due student balance; outcome: first 90-day (60-day) auto delinquency by calendar month \\
Revolving feasibility (Table \ref{tab:cards_probe}) & Ten-percent trade files, archives 2018-07 and 2024-04 & Revolving trades with a positive balance & Formula rate $\alpha = m/B$ from the scheduled minimum and the balance; position from the actual payment against the minimum; duration of the minimum-payment stream at the observed $\alpha$ and a twenty percent annual rate \\
Family counts (Section \ref{sec:autodata}) & One-percent trade files, all archives & Family flags and term ranges of Appendix \ref{appendix:procedures} & Loans, borrowers (distinct consumer keys), first 90-day delinquencies and charge-offs ever (\texttt{results/family\_counts.csv}); student holders at 2023-09 \\
Imputed rates (Appendix \ref{appendix:procedures}) & Origination fields of either panel & Non-mortgage families & Monthly rate by bisection of $m = L r/(1-(1+r)^{-T})$ on $(10^{-7}, 1)$, annualized by twelve, values above $0.6$ discarded \\
Balloon and deferred-payment fields; accommodation codes (Appendix \ref{appendix:strategies}) & One-percent trade files at 2015-01, 2020-01 and 2024-01 (balloons); ten-percent trade files 2019-01 to 2023-10 (accommodation codes) & Open trades by family & Balloon: a positive balloon amount, with and without a due date; accommodation: a forbearance or deferment narrative code on an open trade (\texttt{results/balloon\_probe.csv}, \texttt{forbearance\_probe.csv}) \\
\end{longtable}
\endgroup

\subsection{Data issues: what the files do, what the paper does, where it matters}\label{appendix:data:issues}

Each item states the feature of the raw files, the paper's treatment, and the exhibits it bears on; the section of \citet{gibbs_consumer_2025} that covers the feature is given where one exists.

\paragraph{Files, frequencies and coverage.}

\begin{enumerate}

\item \emph{Two panels at two frequencies.} The one-percent trade files are semiannual from 2005 to 2009 and monthly from January 2010 to December 2025; the ten-percent trade files are 57 snapshots at irregular spacing, semiannual to 2017, quarterly from 2018 to 2022, monthly for part of 2023 and quarterly to April 2024. Every design that needs the timing of an event within a year (the run-off by loan age, the resets, the resumption, the equity margin) runs on the one-percent files; the within-borrower and deficiency designs, the distress transitions and the revolving probe run on the ten-percent files, where the number of loans matters and the timing does not. Persistence is annualized by the actual gap between consecutive ten-percent archives. The survey's Table 4 and Online Appendix G.1 describe the one-percent panel as semiannual from 2005 to 2008 and monthly thereafter; the files on hand are semiannual through 2009 and monthly from 2010.

\item \emph{Scores and attributes.} The one-percent score file is monthly for the trade files' months and the attribute file quarterly; in the one-percent family panels a score is matched at the month of first sight for 99.8 to 99.9 percent of loans and an income at the nearest earlier quarter within a year for 98.3 to 99.4 percent. The ten-percent score file is monthly from 2010, 195 files of about 31 million consumers each, with two score columns both populated; the paper's score is one of the two columns, with the other used on the ten-percent extract where it is missing or zero. The trade files have no age or birth field; age and modeled income exist only in the one-percent attribute file, so on the ten-percent extract they exist only for the one-percent consumers nested in it. Income is the bureau's modeled estimate, not a reported income, as the survey's Section 3.4.4 and Online Appendix E.1 state for every panel; the score is a 24-month ninety-day-delinquency logit built without income (Sections 4.1 and 4.2), which is why score and income enter the paper's fixed effects as separate axes. Matching a score file to a trade extract by consumer key and archive month has one silent failure mode: a key vector that is de-duplicated and then split by a grouping factor of the original length assigns keys to the wrong months and matches a few percent of the extract without an error. The match sets are built by splitting the full key vector by month and de-duplicating within month.

\item \emph{The sampling rule and the key.} The one-percent consumers' keys all end in the same two digits, which are the sampling digits of the panel (the survey's Section 5.2.1 and Online Appendix G.2 describe sampling on the last digits of the Social Security number or of the internal identifier). A subsample defined by the last digits of the key is therefore degenerate.

\item \emph{Archive months differ in coverage.} Some ten-percent archive months hold fewer part files and rows than their neighbours; the October 2018 month holds about 60 percent of a typical month's auto trades. The distress transitions report the absent share of the next archive rather than pooling across months, and every time-series estimate of a hazard runs on the one-percent monthly files.

\item \emph{Origination and first sight.} A loan appears in the files with a lag after its opening date (the survey's Section 3.5.1 and Online Appendix B). Origination amount, payment, term and date are the contract's fields frozen at first sight, not the first-sight balance and archive month, as the survey recommends; cohorts are by the opening date; covariates are matched at the first-sight archive. Loans opened before 2005 enter the one-percent panels at first sight, so the equity margin's 2000 to 2004 cohorts are at risk from 2005.

\item \emph{Reporting-duration limits.} Delinquent trades are deleted seven years after the first delinquency and closed clean trades ten years after the last activity (the survey's Table 3). Terminal outcomes are taken at the last sight of the trade; the persistence measure uses consecutive archives only; no cohort is followed beyond the contract's term.

\item \emph{Transfers between furnishers and stale trades.} A trade transferred to a new servicer stops reporting under the old furnisher key and reappears under a new one, sometimes after a gap (the survey's Online Appendix B and H.2.12). The loan record is keyed by consumer, furnisher and opening date, so a transfer closes one record and opens another with the contract's original opening date. The early-payoff definition counts a cessation of reporting more than six months before maturity as an exit, so the exit hazard includes transfers; a trade that a furnisher stops updating keeps its last rating until it leaves the file. The paper links transfers by matching, within a consumer, records of the same family with the same opening date under different furnisher keys: of the auto exits so classified, 15 percent are transfers, 18 percent in the loan's first year and 17 percent in its second, and the remainder are payoffs, charge-offs and disappearances. The exit hazard the formula uses is net of transfers, by group (\texttt{results/exit\_hazards\_net\_auto.csv}, \texttt{transfer\_links\_auto.csv}; replication \texttt{Code/extract/E18\_transfer\_links.R}).

\end{enumerate}

\paragraph{Populations.}

\begin{enumerate}[resume]

\item \emph{Deceased borrowers.} Furnishers and the bureau lack timely death information, and a trade continues to report after the borrower's death (the survey's Section 3.2.1). The paper applies no age restriction: the unit is the loan, the outcome is the trade's rating, and survival is measured from every exit, so a death that ends payments enters the default or the exit hazard as the trade records it. The one-percent attribute file contains age, and Table \ref{tab:loan_summary} reports its distribution; the survey's recommended removal of consumers with implausible ages would remove trades from the risk set on which the exit hazard is measured.

\item \emph{Fragment records.} One person can hold several credit records (the survey's Section 3.2.1). Consumer keys are used as given. The within-borrower design requires two loans under one key, so a borrower whose loans sit on two fragments is not in that sample; the borrower counts (\texttt{results/family\_counts.csv}) are counts of keys.

\item \emph{Joint and co-signed accounts.} The same trade appears on each liable consumer's record (the survey's Section 3.2.3). The paper reports no population aggregate, so no de-duplication is applied; a joint loan is one record per liable consumer in the family counts and in the loan-level regressions, and in the within-borrower design it is one of each holder's loans.

\item \emph{Non-furnishing lenders.} Some subprime auto lenders, in particular buy-here-pay-here dealers, and most payday and buy-now-pay-later lenders do not furnish (the survey's Section 3.2.2 and Online Appendix D.3). The auto family understates the bottom of the score distribution; the paper states this as coverage and does not correct for it.

\end{enumerate}

\paragraph{Default and status measures.}

\begin{enumerate}[resume]

\item \emph{Charge-off and the manner-of-payment rating.} The charge-off amount is the accounting event and the survey's Online Appendix E.2 notes that its reporting is inconsistent across lenders; the rating (1 current, 2 and 3 for 30 and 60 days past due, 4 for 90, 5 for 120 or more, 8 repossession or foreclosure, 9 charge-off) is populated on every row. On autos the paper reports both: charge-off, the terminal outcome the model prices, and the first rating of 4 or worse, which the survey's Section 3.4.3 recommends as the measure scoring models target. On mortgages charge-off is rare, 10,875 of 6,124,957 loans, the survey's Online Appendix E.2 noting that some lenders never update the status to charge-off and that a trade in a severely derogatory state may leave the file, so the mortgage and home-equity default event is the first 90-day delinquency throughout. The survey's recommended 30-day measure is not used as an outcome because it is the entry to the distress state whose persistence Appendix \ref{appendix:measuredrho} measures.

\item \emph{No payment-history array; persistence from consecutive archives.} The trade files contain the current rating, the date of major delinquency, the last-activity and closure dates and narrative codes, but not the 24- to 84-month payment grid the survey's Section 3.4.3 and Online Appendix C.2 describe. Delinquency dynamics are therefore measured across consecutive archives: on the ten-percent files, among 26.6 million auto trades 30 days past due or worse with no charge-off, $72.7$ percent are still delinquent at the next archive, $6.8$ percent charged off and $20.5$ percent current, a per-gap persistence of $0.80$ and $0.52$ annualized at the mean gap of $4.2$ months; on mortgages the shares are $75.9$, $0.3$ and $23.8$ percent. On the one-percent monthly files the event months of the first 30-, 60-, 90- and 120-day ratings, the first charge-off and the first zero balance are recorded per loan, which dates every event to the month without a grid.

\item \emph{Pandemic accommodations, March 2020 to August 2023.} Under the CARES Act a trade with an accommodation is reported as current, or at its pre-accommodation status, while the accommodation is in effect (the survey's Section 3.4.3, Table 2 and Online Appendix E.2). In the ten-percent files the share of open first-mortgage trades with a forbearance or deferment narrative code rises from $0.04$ to $0.06$ percent in 2019 to $1.3$ percent in July 2020 and returns below $0.1$ percent by October 2021, and the share of open mortgage trades 90 days past due falls from $1.5$ percent in January 2020 to $0.9$ percent in July 2021; on auto trades the flagged share stays below $0.002$ percent and the 90-day share moves from $8.9$ to $8.6$ percent over the same months (\texttt{results/forbearance\_probe.csv}). The narrative code that marks a mortgage forbearance appears nine times as often in 2020 as before (\texttt{results/narrcode\_covid\_candidates.csv}). The paper does not recode accommodated trades as delinquent. The auto estimates are reported by origination cohort, and the cohorts whose windows fall inside the accommodation period, 2019 to 2021, give an elasticity of $0.38$ against $0.40$ for the cohorts before (Section \ref{sec:liquidity}); the equity margin and the resets end in 2018; the resumption's exposure is defined by the end of the federal payment pause.

\item \emph{Student loans: the reporting rule and the on-ramp.} Federal student loans cannot be reported as delinquent before 90 days past due; from March 2020 through September 2023 federally held loans were reported with a zero scheduled payment and no new delinquencies; and the Department of Education's twelve-month on-ramp kept missed payments off the files for a further year, to October 2024 (the survey's Online Appendix D.4). The paper's exposure is the resumed scheduled payment at the end of the pause, and the measured consequence of the on-ramp is in Table \ref{tab:student_takeup}: the share of holders with a past-due student balance stays between two and four percent through January 2025 while about half of the resumed bills go unpaid. Appendix \ref{sec:studentloan} assesses the design under that rule.

\item \emph{Right-censoring and loan age at the archive.} A later archive mechanically shows fewer defaults among young cohorts. Every window requires the contract to be long enough and old enough at the last archive to complete it, and the run-off holds loan age at the archive fixed across bins; a hazard computed across calendar time without that control moves with the age composition of the loans at risk.

\item \emph{The default cost differs by family.} A mortgage charge-off nets the collateral's recovery and a personal-loan charge-off is near-total, so a statistic built on the charged-off amount is not comparable across families. The estimator used here uses the response to the payment and to the horizon and does not use the amount; the balance-scaled penalty channel is measured on autos from the deficiency contrast and on mortgages from the recourse contrast (Section \ref{sec:equity}).

\end{enumerate}

\paragraph{Payments, rates and contract fields.}

\begin{enumerate}[resume]

\item \emph{The mortgage scheduled payment includes escrow.} For mortgages the scheduled payment may include taxes, insurance and association fees, so annuity inversion of the payment misstates the rate (the survey's Section 3.5.4 and Online Appendix D.1). Rate imputation from the payment is restricted to the non-mortgage families, as the survey recommends. On mortgages the paper recovers the note rate from the balance path, the identity in Table \ref{tab:data_four}, which is the survey's balance-based alternative made operational at the month: a rate reset moves the implied rate and leaves the escrow unchanged, an escrow re-analysis moves the escrow and leaves the implied rate within two tenths of a point, and the two are separated on 429,078 payment changes (Table \ref{tab:arm_classes}, Appendix \ref{sec:armreset}). The mortgage payment margin of Table \ref{tab:equity_margin} is estimated on the reported payment, escrow included.

\item \emph{Garbage payment rows.} The scheduled payment is in the millions on a small share of rows and the product of payment and term overflows a 32-bit integer. The panel estimates keep origination payments between \$20 and \$100,000 and convert to double precision before any product.

\item \emph{Balloon and deferred-payment fields.} The one-percent trade files contain a balloon amount and due date and a deferred-payment start date. Trades with a positive balloon amount number 2,144 to 2,671 of 1.3 to 1.4 million open auto trades and, with a due date, 2,389 to 4,578 of 0.9 to 1.5 million open first mortgages in the archives of January 2015, 2020 and 2024 (\texttt{results/balloon\_probe.csv}); Appendix \ref{appendix:strategies} records the design they would support and the events they hold.

\end{enumerate}

\paragraph{Silent failure modes in the files.} Three features of the raw files produce a wrong answer without an error, and each is handled when the file is read. Date fields are variable-width integers in month, day, year order (8012002 is 1 August 2002), a day of zero occurs, and 1.6 percent of opening dates do not parse; dates are zero-padded to eight digits before parsing, and rows whose opening date fails to parse are dropped from the timing calculations only. The consumer and furnisher keys span the full 64-bit integer range, so parsing the keys as numeric loses the low digits of nineteen-digit keys and merges consumers; keys are read as character strings in every script. A column with no values in one part file is typed numeric there and character in another, so a join across part files matches nothing; date and key columns are cast to character immediately after each file is read.

\paragraph{Where the paper's practice differs from the survey's guidance.} Four practices differ, each for a reason of the design. The paper applies no age restriction and removes no deceased consumers, because the risk set for the exit hazard is every trade (item 8). It does not link trades across furnisher transfers, because the record is keyed by furnisher so that origination fields freeze once, at the cost of counting transfers in the exit hazard (item 7). It does not recode accommodated trades as delinquent during the pandemic, and reports the affected cohorts separately instead (item 14). And its headline auto outcome is charge-off alongside the survey's 90-day measure, because the model prices the terminal outcome (item 12). The remaining guidance is followed: origination fields rather than first-sight balances; tradeline-level rather than aggregated measures; the annuity inversion for installment rates and its exclusion for mortgages; and the four disclosures of Table \ref{tab:data_four}.

\subsection{What this panel adds to the survey's account}\label{appendix:data:adds}

The files on hand differ from the survey's account of this panel, and of tradeline extracts in general, in the following respects. The one-percent panel is monthly from January 2010 through December 2025 and semiannual from 2005 to 2009, where the survey's Table 4 and Online Appendix G.1 date the monthly files from 2009; the ten-percent trade files are 57 archives at irregular spacing against a monthly score file from 2010 with two populated score columns and keys matching the trade files in full, where the survey's Online Appendix G.1 describes the ten-percent product as monthly; and attributes and trade files are held for different samples, age and income for the one-percent consumers only, so that a covariate-matched subsample of the ten-percent extract is the nested one-percent panel's loans. The consumer key contains the sampling digits, so subsamples on its last digits are degenerate (item 3), and matching scores to a trade extract by key and month has the failure mode of item 2, whose check is the matched share against the score file's coverage. On the measures: the amortization identity on the balance path separates a rate reset from an escrow re-analysis at monthly frequency on every mortgage without a contract flag, one payment change in ten at the hybrid reset ages on 2004 to 2008 first mortgages being a rate change and the median escrow re-analysis ninety dollars (item 18); on long contracts the charge-off amount is unusable as the default event at any frequency, 10,875 charge-offs on 6.1 million first mortgages, while the rating is populated on every row (item 12); without the payment-history grid, consecutive-archive transitions measure the persistence of distress, $0.80$ per $4.2$-month gap and $0.52$ annualized with one fifth of delinquent trades curing per gap, and their absent share of 14 to 19 percent is a component a grid-based measure does not report (item 13); and loan age at the archive is the censoring control for any time-series estimate of a default-based statistic, the run-off by loan-age bin among survivors being the form in which the payment margin can be compared across calendar time (item 16). On the period: the student on-ramp's reporting rule keeps the past-due share of holders between two and four percent through January 2025 while about half of the resumed bills go unpaid (item 15); the pandemic accommodations are visible on mortgages, a forbearance code on $1.3$ percent of open first mortgages at the peak and a fall in the 90-day share by two fifths, and not on auto trades, where the flagged share is below $0.002$ percent (item 14); and balloon and deferred-payment fields exist and are rare, about two thousand open auto trades and two to five thousand dated mortgage balloons per one-percent archive (item 20). The three silent failure modes of the preceding subsection, packed dates with a zero day, keys beyond double precision and per-file typing of all-missing columns, complete the list.
